\section{结论}

本章针对常值延迟问题提出了一种基于模型的方法。
核心思想是学习信念的向量表征，将无限维量转化为有限维。
这得益于对\gls{nn}的精心选择。
该信念表征可作为状态的预处理模块接入任何\gls{rl}算法，正如本文将其接入\gls{trpo}构建 D-TRPO 一样。
通过一个简单技巧，展示了如何利用 D-TRPO 在较大延迟下学习的策略，直接应用于更小延迟。
该技巧显著降低了适应延迟的学习成本，因为可以同时学习多个延迟。
实验评估证实了这一可预见的结果，策略在任何更小延迟下均取得相似的性能。

随后对常值延迟问题进行了分析。
延迟越长、最优性能越低这一事实已得到形式化证明，该证明可为未来分析提供有用工具。
此外，揭示了基于模型方法的复杂度结果，展示了该方法的局限性。
基于相同思路，给出了一个反例，证明与最优延迟策略相比，基于模型的策略在某些\glspl{mdp}中可能产生次优的期望回报或平均奖励。

最后，通过实验分析评估了 D-TRPO 的能力。
实验表明，尽管 A-SAC 是整体上最有效的方法，但 D-TRPO 能够适应广泛场景，特别是随机\glspl{mdp}，甚至也包括确定性环境。
因此，信念表征网络是一种通用方法，可作为任何\gls{rl}算法的即插即用预处理模块。
值得注意的是，它允许控制信念表征的大小，从而控制\gls{rl}算法输入的维度。

一个有价值的未来方向是利用模型的知识将其融入\gls{rl}算法内部，例如增强演员-评论家方法中的评论家更新。
